\documentclass[12pt]{article}
\usepackage{amsmath, amssymb}
\usepackage[margin=1in]{geometry}
\setlength{\parskip}{6pt}
\title{QUBO Formulation: Energy Storage and Microgrid Siting Problem}
\date{}
\begin{document}
\maketitle

\subsection*{Energy Storage and Microgrid Siting Problem}

\paragraph{Problem Statement.}
Given a set of $N$ candidate sites for energy storage systems or microgrids, each site $i \in \{0,\dots,N-1\}$ is characterized by a resilience benefit score $B_i$, an investment cost $C_i$, and a pairwise overlap penalty $P_{ij}$ for every pair of distinct sites $i,j$. The objective is to select exactly $K$ sites such that the total siting score is maximized. The total siting score is defined as the sum of the resilience benefits of the selected sites, minus the sum of their investment costs, minus the sum of pairwise overlap penalties between all selected site pairs.

\paragraph{Decision Variables.}
For each candidate site $i \in \{0,\dots,N-1\}$, introduce a binary decision variable $x_i \in \{0,1\}$. The variable $x_i$ equals $1$ if site $i$ is selected for deployment, and $0$ otherwise.

\paragraph{QUBO Derivation.}
Step 1 --- The original maximization objective is $\max \sum_{i=0}^{N-1} B_i x_i - \sum_{i=0}^{N-1} C_i x_i - \sum_{0 \le i < j \le N-1} P_{ij} x_i x_j$. Converting to a minimization problem for QUBO yields
\begin{align}
  \min \sum_{i=0}^{N-1} (C_i - B_i) x_i + \sum_{0 \le i < j \le N-1} P_{ij} x_i x_j.
\end{align}

Step 2 --- The exact-$K$ selection constraint is
\begin{align}
  \sum_{i=0}^{N-1} x_i = K.
\end{align}

Step 3 --- We enforce this constraint via a quadratic penalty term with weight $\lambda > 0$: $H_{\text{pen}}(x) = \lambda \left( \sum_{i=0}^{N-1} x_i - K \right)^2$. Expanding the square and applying the idempotent property $x_i^2 = x_i$ for binary variables gives
\begin{align}
  H_{\text{pen}}(x) &= \lambda \left( \sum_{i=0}^{N-1} x_i^2 + \sum_{i \neq j} x_i x_j - 2K \sum_{i=0}^{N-1} x_i + K^2 \right) \\
  &= \lambda \left( \sum_{i=0}^{N-1} x_i + 2 \sum_{0 \le i < j \le N-1} x_i x_j - 2K \sum_{i=0}^{N-1} x_i + K^2 \right) \\
  &= \sum_{i=0}^{N-1} \lambda(1 - 2K) x_i + 2\lambda \sum_{0 \le i < j \le N-1} x_i x_j + \lambda K^2.
\end{align}

Step 4 --- Combining the objective and penalty terms yields the total QUBO Hamiltonian $H(x)$:
\begin{align}
  H(x) &= \sum_{i=0}^{N-1} (C_i - B_i) x_i + \sum_{0 \le i < j \le N-1} P_{ij} x_i x_j + \sum_{i=0}^{N-1} \lambda(1 - 2K) x_i + 2\lambda \sum_{0 \le i < j \le N-1} x_i x_j + \lambda K^2 \\
  &= \sum_{i=0}^{N-1} \left[ C_i - B_i + \lambda(1 - 2K) \right] x_i + \sum_{0 \le i < j \le N-1} \left[ P_{ij} + 2\lambda \right] x_i x_j + \lambda K^2.
\end{align}

Step 5 --- Reading off the coefficients from the standard QUBO form $H(x) = \sum_{i} Q_{i,i} x_i + \sum_{i < j} Q_{i,j} x_i x_j + c$, we obtain the general closed-form expressions for the upper-triangular entries of the QUBO matrix $Q$ and the constant offset $c$:
\begin{align}
  Q_{i,i} &= C_i - B_i + \lambda(1 - 2K), \quad \forall i \in \{0,\dots,N-1\}, \\
  Q_{i,j} &= P_{ij} + 2\lambda, \quad \forall 0 \le i < j \le N-1, \\
  c &= \lambda K^2.
\end{align}

\paragraph{Penalty Weight Justification.}
The penalty weight $\lambda$ must be chosen sufficiently large to ensure that any violation of the exact-$K$ constraint results in a strictly higher energy than any feasible solution. The maximum possible gain in the original objective from altering the number of selected sites by one is bounded by $\max_{i} |B_i - C_i|$. Therefore, setting $\lambda > \max_{i} |B_i - C_i|$ guarantees that the quadratic penalty incurred by selecting $K \pm 1$ sites outweighs any potential benefit or cost saving from the linear and quadratic objective terms, strictly enforcing feasibility at the global minimum.

\paragraph{Correctness Argument.}
Minimizing $H(x)$ over $\{0,1\}^N$ recovers the optimal site selection because (i) all feasible solutions satisfying $\sum_{i=0}^{N-1} x_i = K$ incur zero penalty and thus evaluate exactly to the original objective function, and (ii) any infeasible solution violates the constraint by at least one unit, incurring a penalty increase of at least $\lambda$. Since $\lambda$ exceeds the maximum possible objective improvement from any single-site modification, every infeasible configuration has an energy strictly greater than the optimal feasible configuration. Consequently, the global minimum of $H(x)$ must correspond to a feasible assignment that maximizes the original siting score.

\paragraph{Illustrative Example.}
Consider a concrete instance with $N=3$ candidate sites, $K=2$ required selections, and penalty weight $\lambda=10$. Let the benefits be $B_0=10, B_1=8, B_2=12$, costs be $C_0=5, C_1=4, C_2=6$, and overlap penalties be $P_{01}=2, P_{02}=3, P_{12}=1$. Note that $\max_{i} |B_i - C_i| = 6 < \lambda$, satisfying the penalty condition. Substituting these values into the general formulas yields the following concrete QUBO coefficients:
\begin{align}
  Q_{0,0} &= 5 - 10 + 10(1 - 4) = -35, \\
  Q_{1,1} &= 4 - 8 + 10(1 - 4) = -34, \\
  Q_{2,2} &= 6 - 12 + 10(1 - 4) = -36, \\
  Q_{0,1} &= 2 + 20 = 22, \quad Q_{0,2} = 3 + 20 = 23, \quad Q_{1,2} = 1 + 20 = 21, \\
  c &= 10 \cdot 2^2 = 40.
\end{align}
The resulting Hamiltonian is $H(x) = -35x_0 - 34x_1 - 36x_2 + 22x_0x_1 + 23x_0x_2 + 21x_1x_2 + 40$. Evaluating $H(x)$ over the three feasible assignments ($x_0+x_1+x_2=2$) gives $H(1,1,0)=-7$, $H(1,0,1)=-8$, and $H(0,1,1)=-9$. The global minimum is $-9$ at $x^*=(0,1,1)$, indicating that sites 1 and 2 should be deployed to achieve the maximum total siting score.

\end{document}